\section{Preparing local tests}
\label{sec:local-tests}

Fix $n\geq1$ and a prime $p$, and suppose that $\Zp$ acts faithfully
and jointly continuously on a connected Hausdorff second-countable
topological $n$-manifold $M$ without boundary, as in \Cref{thm:main}.
We will construct a degree-one test from a compact sphere through an
orbit space, and then move a positive-degree multiple of that test
into a prescribed open region of the orbit space.  The construction
begins with a faithful action inside one coordinate chart.

\subsection{An invariant chart subset and a compact test source}

\begin{proposition}[An invariant chart subset]
\label{prop:chart-reduction}
There are an open subgroup $G\cong\Zp$ and a nonempty connected
$G$-invariant open subset $O$ of a coordinate chart of $M$ such that
$G$ acts faithfully and preserves orientation on $O$.  In particular,
we may regard $O$ as the actual open subset of $\R^n$ supplied by that
chart.
\end{proposition}

\begin{proof}
We use Newman's theorem in the following form: a finite-order
homeomorphism of a connected topological manifold which is the
identity on a nonempty open subset is the identity everywhere
\cite{Newman1931}; for a modern statement and proof, see
\cite[author manuscript, pp.~2--3]{Pardon2019}. The local reduction using this
theorem also appears in \cite[\S4.1]{Pardon2013}.

There is a point $x\in M$ such that no open subgroup of $\Zp$ fixes a
neighborhood of $x$ pointwise.  To prove this, suppose otherwise, and
write
\[
 H_j=p^j\Zp,
 \qquad F_j=\operatorname{int}_M\operatorname{Fix}(H_j).
\]
Some $F_j$ is nonempty.  We show that it is closed.  If
$x\in\overline{F_j}$, then $H_j$ fixes $x$, since its fixed set is
closed.  By the supposition, some neighborhood $N$ of $x$ is fixed
pointwise by $H_\ell$.  Compactness of $H_j$ and continuity at
$H_j\times\{x\}$ give a coordinate ball $B\ni x$ so small that
$H_jB\subset N$.  The saturation $N'=H_jB$ is open and connected:
every translate of $B$ contains $x$.  It is $H_j$-invariant, and the
restricted action factors through the finite group
$H_j/(H_j\cap H_\ell)$.  Each of its homeomorphisms is therefore
periodic and fixes the nonempty open subset $N'\cap F_j$ pointwise.
Newman's theorem makes each the identity on $N'$.  Thus $x\in F_j$,
proving closedness.  Connectedness gives $F_j=M$, contrary to
faithfulness.

Choose such an $x$, a chart $C\ni x$, and a coordinate ball $B\ni x$
with $\overline B$ compactly contained in $C$.  Joint continuity and
compactness of $\overline B$ give an open subgroup $H$ such that
$H\overline B\subset C$.  Shrinking $H$ further gives $Hx\subset B$.
Then
\[
 O=HB=\bigcup_{g\in H}gB\subset C
\]
is open and invariant, and is connected because every $gB$ meets
$B$ at $gx$.  Its restriction kernel in $H$ is closed.  A nonzero
closed subgroup of $\Zp$ is $p^j\Zp$ for some $j$: choose an element
of smallest finite $p$-adic valuation and use the density of its
integer multiples in that subgroup.  Thus a nonzero restriction
kernel would be an open subgroup fixing the neighborhood $O$ of
$x$, which is impossible.

The chart orients $O$.  Every homeomorphism of connected $O$ either
preserves or reverses this orientation, so every square preserves
it.  Replace $H$ by $G=2H$.  This is an open subgroup isomorphic to
$\Zp$ and its restricted action remains faithful.  For odd $p$ this
replacement leaves $H$ unchanged; for $p=2$ it has index two.
\end{proof}

Fix an odd integer $d>n$ with $d\geq3$ for the rest of the proof.  Give
\[
 W=O\times\R^{d-n}\subset\R^d
\]
the product orientation, and let $G$ act trivially on the auxiliary
factor.  This action remains faithful and preserves orientation.
Write $W^+=W\cup\{\infty\}$ for its one-point
compactification and use
$S^d=\R^d\cup\{\infty_{\mathrm{amb}}\}$ with the matching
orientation as the compact test source.  Its orientation form represents
$u_d$ of \Cref{thm:sphere-lattice}.  The choice of $d$ will allow
us to apply the odd-sphere homotopy calculation below.

\begin{lemma}[Collapse and compactification]
\label{lem:collapse}
The map
\[
 c:S^d\longrightarrow W^+,
 \qquad
 c(z)=
 \begin{cases}
 z,&z\in W,\\
 \infty,&z\notin W,
 \end{cases}
\]
is continuous.  The action on $W$ extends to a jointly continuous
action on $W^+$ fixing $\infty$.  Both $W^+$ and its orbit space by
any closed subgroup of $G$ are compact metrizable spaces.
\end{lemma}

\begin{proof}
The inverse image of an open subset of $W$ is open in $S^d$.  A
basic neighborhood of $\infty$ is $W^+\setminus K$, where
$K\subset W$ is compact; its inverse image is the open set
$S^d\setminus K$.  This proves continuity of $c$ for an arbitrary
open $W$.

For compact $K\subset W$, its saturation $GK$ is compact, being the
image of $G\times K$.  Every $g\in G$ maps
$W^+\setminus GK$ into $W^+\setminus K$.  This gives joint
continuity at all $(g,\infty)$, while continuity on $G\times W$ was
already established.

The space $W$ is locally compact, noncompact, and second countable;
its one-point compactification is compact metrizable.  If $d_0$ is
a compatible metric on $W^+$, then
\[
 d_1(x,y)=\max_{g\in G}d_0(gx,gy)
\]
is a compatible invariant metric, by uniform continuity on the
compact space $G\times W^+$.  For a closed subgroup $H\subset G$,
the metric
\[
 d_H(Hx,Hy)=\min_{h\in H}d_1(x,hy)
\]
induces the quotient topology.  Its positivity on distinct orbits
follows from compactness, and the triangle inequality follows from
invariance.  This proves the last assertion.
\end{proof}

\begin{lemma}[An averaged degree-one test]
\label{lem:averaged-test}
After replacing $G$ by an open subgroup and identifying that
subgroup with $\Zp$, set
\[
 q:W^+\longrightarrow Y=W^+/G,
 \qquad y_\infty=q(\infty).
\]
There is a continuous map $a:Y\to S^d$ such that
\begin{equation}
 \deg(aqc)=1.
 \label{top:degree-one}
\end{equation}
Moreover, $aq$ is homotopic on $W^+$ to a degree-one pinch supported
in a round ball compactly contained in $W$.
\end{lemma}

\begin{proof}
Fix $\xi\in S^d$ and a map $P:\R^d\to S^d$ which is constant
with value $\xi$ outside the unit ball and induces an
orientation-preserving homeomorphism
$\overline B^d/\partial B^d\to S^d$.  Choose $z\in W$ and $r>0$
with $\overline B(z,r)\subset W$.  The map
\[
 P_{z,r}(w)=P((w-z)/r),\qquad P_{z,r}(\infty)=\xi,
\]
is continuous on $W^+$ and $P_{z,r}c$ has degree one.

Regard $S^d$ as the unit sphere in $\R^{d+1}$.  By continuity and
compactness, a sufficiently small open subgroup $H\subset G$
satisfies
\[
 \sup_{h\in H,\,w\in W^+}
 \|P_{z,r}(hw)-P_{z,r}(w)\|<\tfrac12.
\]
Let $\mu_H$ be Haar probability measure and put
\[
 A(w)=\int_H P_{z,r}(hw)\,d\mu_H(h).
\]
This is a continuous $H$-invariant $\R^{d+1}$-valued map and
$\|A(w)-P_{z,r}(w)\|<1/2$.  Thus every vector
$(1-t)P_{z,r}(w)+tA(w)$ has norm at least $1/2$.  Normalizing these
vectors gives a homotopy from $P_{z,r}$ to the invariant map
$A/\|A\|$.  The homotopy fixes $\infty$ and is constant with value
$\xi$ outside the compact subset $H\overline B(z,r)$ of $W$.
Replace $G$ by $H$.  The normalized map descends to $a:Y\to S^d$,
and the homotopy proves \eqref{top:degree-one}.  Faithfulness and
orientation preservation persist under this replacement.
\end{proof}

We henceforth fix $G,q,Y,a$ as in \Cref{lem:averaged-test}.

\subsection{Rational orbit cohomology and maps to an odd sphere}

To move the test, we first detect rational cohomology on the orbit
space by pulling it back to the compactification.  Rational
cohomology of compact metrizable spaces below means \v Cech
cohomology, computed from finite open-cover nerves. We use its natural
agreement with cohomology of the constant sheaf; see
\cite[\S1.6 and \S2]{BergfalkLupiniPanagiotopoulos2024}.
The sheaf \v Cech comparison is an isomorphism on compact Hausdorff
spaces and commutes with pullback
\cite[Lemmas~20.16.2 and~20.15.1]{StacksCechComparison}.
To relate it to nerve cochains, shrink a finite cover so that the
closure of each member lies in its original member. On each compact
intersection of these closures, any locally constant cochain value has
finite image and a positive constancy scale. For finitely many
cochains, a sufficiently fine common finite refinement therefore
makes every pulled-back value constant on each new intersection.
Applying this to cocycles and to their primitives identifies the two
direct limits. The inclusions of constant cochains into locally
constant cochains commute with refinement and pullback, so the
identification is natural. No covering-dimension bound is used.

\begin{lemma}[Rational orbit pullback]
\label{lem:orbit-injectivity}
Let $E$ be a compact metrizable space with a continuous $G$-action,
and let $r:E\to E/G$ be its orbit map.  For any field $k$, in the
bounded-below derived category of sheaves one has
\[
 Rr_*\const{k}{E}\simeq r_*\const{k}{E}.
\]
For $k=\Q$ the unit
$\const{\Q}{E/G}\to r_*\const{\Q}{E}$ has a natural sheaf
retraction.  Consequently
\[
 r^*:\check H^j(E/G;\Q)\longrightarrow\check H^j(E;\Q)
\]
is injective for every $j$.
\end{lemma}

\begin{proof}
Every orbit is homeomorphic to $G/G_x$ for a closed subgroup $G_x$.
Such a quotient is a compact zero-dimensional space: for $\Zp$ it
is either a point, a finite quotient, or $\Zp$ itself.  On a compact
zero-dimensional space, global sections of sheaves of vector
spaces form an exact functor.  Indeed, local lifts through a sheaf
epimorphism can be chosen on a finite cover and then glued over a
finite disjoint clopen refinement.  Thus the higher cohomology of
each orbit vanishes.

The map $r$ is closed and separated and has compact fibers.
Bounded-below proper base change
\cite[Lemma~20.18.1 and Theorem~20.18.2]{StacksProperBaseChange}
therefore gives
\[
 (R^jr_*\const{k}{E})_y
   =H^j(r^{-1}(y);\const{k}{r^{-1}(y)})=0
   \quad(j>0).
\]
This use of proper base change imposes no covering-dimension bound
on $E$ or $E/G$.

For an open $U\subset E/G$ and a locally constant function
$\varphi:r^{-1}(U)\to\Q$, define
\begin{equation}
 \mathcal A_U(\varphi)(r(x))
     =\int_G\varphi(gx)\,d\mu_G(g).
 \label{top:rational-average}
\end{equation}
The expression is independent of the representative $x$.  For
fixed $x$, the locally constant function $g\mapsto\varphi(gx)$ on
$G$ factors through some finite quotient $G/p^\ell G$, so its average
is rational.

It is also locally constant as a function of $x$.  For every
$g\in G$, continuity gives a product neighborhood of $(g,x)$ on
which $\varphi(hx')$ is constant.  A finite collection of the group
neighborhoods covers $G$.  Intersecting the corresponding
neighborhoods of $x$ shows that
$\varphi(hx')=\varphi(hx)$ for all $h\in G$ and all $x'$ in one
neighborhood of $x$.  Their averages are equal there.  An invariant
locally constant function descends to a locally constant function
on the orbit space, because $r$ is an open quotient map.
The operators \eqref{top:rational-average} commute with
restriction and fix functions pulled back from $U$.  They give
the asserted sheaf retraction.  Applying derived global sections,
and using $Rr_*\const{\Q}{E}=r_*\const{\Q}{E}$, makes the usual
cohomology pullback split injective.  The natural \v Cech--sheaf
comparison gives the stated conclusion.
\end{proof}

Let $\gamma_d\in\check H^d(S^d;\Q)$ denote the rational orientation
class for our chosen sphere orientation.  On an odd sphere, rational
triviality can be strengthened to nullhomotopy after a positive-degree
postcomposition.  The next proof passes through a finite nerve, so
it applies to a compact metrizable domain of arbitrary dimension.
The finite-cover construction is classical \v Cech theory
\cite[Chapter~IX]{EilenbergSteenrod1952}; see also
\cite[\S2 and \S4.3]{BergfalkLupiniPanagiotopoulos2024}.
We construct the sphere-valued approximation explicitly.

\begin{lemma}[Killing a rationally trivial sphere map]
\label{lem:odd-sphere-null}
Let $d\geq3$ be odd.  If $K$ is compact metrizable and $v:K\to S^d$ satisfies
$v^*\gamma_d=0$, then there is a positive integer $s$ and a degree
$s$ map $D_s:S^d\to S^d$ for which $D_sv$ is nullhomotopic.
\end{lemma}

\begin{proof}
The case $K=\varnothing$ is immediate.  Choose a finite open cover
$\mathcal U=\{U_i\}$ of $K$ on each of whose members $v$ has image
of Euclidean diameter less than $1/4$, and choose $x_i\in U_i$.
For every simplex $\sigma$ of the nerve $N\mathcal U$, choose a
point in the corresponding intersection.  All the vectors
$v(x_i)$ for vertices of $\sigma$ are within $1/4$ of its image.
Their convex combinations are therefore nonzero.  The formula
\[
 b\left(\sum_i t_i i\right)
    =\frac{\sum_i t_i v(x_i)}{\|\sum_i t_i v(x_i)\|}
\]
defines a continuous map $b:|N\mathcal U|\to S^d$.
A subordinate partition of unity defines a nerve map
$p_{\mathcal U}:K\to|N\mathcal U|$.  Normalized interpolation
between $v(x)$ and the corresponding convex combination gives
$v\simeq bp_{\mathcal U}$.

For finite covers of a compact space, \v Cech cohomology is the
direct limit of the cohomologies of their nerves.  The map to the
limit from $H^d(|N\mathcal U|;\Q)$ is the pullback by
$p_{\mathcal U}$; this follows from the vertex-star cover and is
independent of the subordinate partition.  Since the image of
$b^*\gamma_d$ is zero, some finite refinement $\mathcal V$ and its
simplicial refinement map
$\rho:|N\mathcal V|\to|N\mathcal U|$ satisfy
\[
 (b\rho)^*\gamma_d=0\quad\hbox{in }H^d(|N\mathcal V|;\Q).
\]
Moreover $p_{\mathcal U}\simeq\rho p_{\mathcal V}$: at each
point, the vertices appearing in both barycentric expressions
belong to cover members containing that point, so linear
interpolation stays in one simplex of $N\mathcal U$.  We have
reduced the assertion to the map
$b\rho:P\to S^d$ from the finite complex $P=|N\mathcal V|$.

We now use that $d$ is odd and $d\geq3$.  Using based cellular degree
maps, consider the mapping telescope
\begin{equation}
 \mathcal S=\operatorname{Tel}
 \bigl(S^d\xrightarrow{D_2}S^d
       \xrightarrow{D_3}S^d\xrightarrow{D_4}\cdots\bigr).
 \label{top:odd-telescope}
\end{equation}
We verify its homotopy type directly.  Serre's finiteness theorem
gives finite $\pi_j(S^d)$ for $j>d$
\cite[Chapter~V, \S3, Proposition~3]{Serre1951}.
We also use his precise degree-map assertion, valid in every
homotopy degree of an odd sphere:
if $q x=0$ in a homotopy group, then postcomposition with a
degree-$2q$ map kills $x$
\cite[Chapter~V, \S1, Corollary~2 to Proposition~1]{Serre1953}.
The telescope has CW type, and its homotopy groups are the direct
limits of those of its spheres: maps and homotopies from spheres
have bounded height, and each finite subtelescope retracts onto
its final sphere.  Any element of $\pi_j(S^d)$, $j>d$, has some
finite order $q$.  A sufficiently long subsequent product of the
degrees $2,3,4,\ldots$ is divisible by $2q$; its sphere map is
homotopic to a composite through a degree-$2q$ map, and therefore
kills that element.  Thus all these higher direct limits vanish.
In degree $d$ the direct limit is
$\operatorname{colim}(\Z\xrightarrow{2}\Z
\xrightarrow{3}\Z\to\cdots)=\Q$, and in degrees below $d$
it is zero.  Hence $\mathcal S$ has type $K(\Q,d)$.
The same odd-sphere telescope and finite-stage argument are used in
\cite[proof of Proposition~4.5, p.~928]{DranishnikovFerryWeinberger2003}.

Let $j:S^d\to\mathcal S$ include the first sphere.  Choose
$\gamma_{\mathcal S}\in H^d(\mathcal S;\Q)$ with
$j^*\gamma_{\mathcal S}=\gamma_d$.  Indeed, $j_*$ in degree-$d$
homology is the inclusion $\Z\hookrightarrow\Q$, and universal
coefficients identify this restriction map with the isomorphism
$\Hom(\Q,\Q)\to\Hom(\Z,\Q)$.
Then $(jb\rho)^*\gamma_{\mathcal S}=(b\rho)^*\gamma_d=0$.
Ordinary Eilenberg--MacLane representability on the finite complex $P$
makes $jb\rho$ nullhomotopic.  A nullhomotopy has compact image.  The
continuous height function on the telescope is therefore bounded
on that image, so the image lies in a finite subtelescope.
Retracting this subtelescope onto its last sphere gives a
nullhomotopy of $D_s b\rho$, where $D_s$ is the composite of
finitely many of the displayed degree maps and $s>0$ is their
product.  Composing with $p_{\mathcal V}$ and using
$v\simeq b\rho p_{\mathcal V}$ proves the assertion.
\end{proof}

\subsection{Concentrating the test on finite quotient sheets}

Faithfulness supplies a region with finite quotient sheets for every
quotient order, even when the original action has fixed points or
nontrivial stabilizers.  The degree-one test $a$ remains fixed while
these regions vary.

\begin{lemma}[Finite quotient sheets]
\label{lem:finite-sheets}
For every integer $k\ge1$, put $m=p^k$ and
$\zeta=\exp(2\pi i/m)$.  There exists a nonempty open
$V\subset Y\setminus\{y_\infty\}$ such that the finite quotient
map
\[
 W/(p^kG)\longrightarrow W/G=Y\setminus\{y_\infty\}
\]
over $V$ is a disjoint union of $m$ homeomorphisms onto $V$.
There is a locally constant function
\[
 \beta:q^{-1}(V)\longrightarrow\{1,\zeta,\ldots,\zeta^{m-1}\}
 \subset S^1
\]
satisfying
\begin{equation}
 \beta(gw)=\zeta^{\bar g}\beta(w),
 \qquad \bar g\in G/p^kG=\Z/m.
 \label{top:sheet-multiplier}
\end{equation}
\end{lemma}

\begin{proof}
Some $x\in W$ has stabilizer $G_x\subset p^kG$.  Otherwise,
the classification of closed subgroups of $\Zp$ would imply
$p^{k-1}G\subset G_x$ for every $x\in W$, contradicting
faithfulness.  Set $H=p^kG$, $Z=W/H$, and let $z$ be the image
of $x$.  This is a Hausdorff orbit space, and the stabilizer of
$z$ for the finite action of $G/H$ is
$(G_x+H)/H=0$.  Write $\tau$ for the element induced by $1\in G$.

Choose pairwise disjoint open neighborhoods $U_j$ of
$\tau^jz$, $0\le j<m$, and put
$N=\bigcap_{j=0}^{m-1}\tau^{-j}U_j$.  Then the translates
$\tau^jN$ are pairwise disjoint.  If $r:Z\to W/G$ is the
finite quotient map and $V=r(N)$, openness of $r$ shows that
$V$ is nonempty and open, and
\[
 r^{-1}(V)=\coprod_{j=0}^{m-1}\tau^jN.
\]
Each restriction $r:\tau^jN\to V$ is an open continuous
bijection, hence a homeomorphism.  Pulling these sheets back under
$W\to Z$ partitions $q^{-1}(V)$ into relatively clopen subsets.
Give the pullback of $\tau^jN$ the constant value $\zeta^j$.
The resulting function $\beta$ is locally constant, and the cyclic
permutation of the sheets proves
\eqref{top:sheet-multiplier}.
\end{proof}

To concentrate a positive-degree postcomposition of $a$ in such a
region, we will extend a nullhomotopy from its closed complement.
The following construction provides the required extension.

\begin{lemma}[Extending a homotopy from a closed subset]
\label{top:neighborhood-extension}
Let $Y_0$ be compact metrizable, $F_0\subset Y_0$ closed, and
$f:Y_0\to S^d$ continuous.  Every homotopy
$H:F_0\times[0,1]\to S^d$ starting at $f|_{F_0}$ is the
restriction of a homotopy on $Y_0$ starting at $f$.
\end{lemma}

\begin{proof}
Put $I=[0,1]$.  Paste $H$ with $f$ on the closed subset
\[
 C=(F_0\times I)\cup(Y_0\times\{0\})\subset Y_0\times I.
\]
Viewing the resulting map as $\R^{d+1}$-valued, extend each of its $d+1$
coordinates by the Tietze extension theorem to obtain
$E:Y_0\times I\to\R^{d+1}$.  The open set
$N=\{\|E\|>1/2\}$ contains $C$, and $E/\|E\|$ extends the
pasted sphere map to $N$.  Compactness of $I$ gives an open
$U\supset F_0$ with $U\times I\subset N$; explicitly, take the
complement of the projection of the compact set
$(Y_0\times I)\setminus N$.
Choose a continuous $\lambda:Y_0\to[0,1]$ equal to one on $F_0$
and zero on $Y_0\setminus U$.  Then
\[
 \widetilde H(y,t)
   =\frac{E(y,t\lambda(y))}{\|E(y,t\lambda(y))\|}
\]
is defined everywhere: if $y\in U$ its argument lies in
$U\times I$, and otherwise its argument is $(y,0)\in C$.
It is continuous, starts at $f$, and agrees with $H$ on
$F_0\times I$.
\end{proof}

We can now concentrate a positive-degree postcomposition of the fixed
test $a$ into any prescribed nonempty open region.  The degree $s$
below may depend on that region; for a region supplied by
\Cref{lem:finite-sheets}, this allows dependence on $k$.

\begin{proposition}[Moving a test into a prescribed open set]
\label{prop:move-test}
For every nonempty open $V\subset Y\setminus\{y_\infty\}$, set
$F=Y\setminus V$.  There are a positive integer $s$, a degree
$s$ map $D_s:S^d\to S^d$, and a map $h:Y\to S^d$ such that
\[
 h\simeq D_sa,
 \qquad h|_F\equiv b
\]
for some $b\in S^d$.  In particular $h(y_\infty)=b$ and
$q^{-1}h^{-1}(S^d\setminus\{b\})\subset q^{-1}(V)\subset W$.
\end{proposition}

\begin{proof}
We first show $(a|_F)^*\gamma_d=0$.  Choose a round ball with
closure in the nonempty open set $q^{-1}(V)\subset W$.  The
pinch in \Cref{lem:averaged-test} can be moved to this ball by a
homotopy on $W^+$.  Here is explicit control of that homotopy.
Connect the two centers by a path in connected open $W$, which
is path connected.  Its compact image has positive distance
from $\R^d\setminus W$.  Choose a sufficiently small fixed
radius for transport along the path; at the two ends shrink
and enlarge radially as needed.  Applying the formula
$P_{z(t),r(t)}$ from \Cref{lem:averaged-test} gives a continuous
family of pinches.  All radii are positive and all closed balls
lie in one compact subset of $W$.  Consequently the family is
also continuous at $\infty$, where it has constant value $\xi$.

It follows that $aq$ is homotopic to a map constant on
$q^{-1}(F)$.  Thus the pullback of $(a|_F)^*\gamma_d$ to
$q^{-1}(F)$ is zero.  This preimage is a closed invariant compact
subset of $W^+$, and its orbit space is $F$.  The restricted orbit map
$q^{-1}(F)\to F$ therefore meets the hypotheses of
\Cref{lem:orbit-injectivity}, so $(a|_F)^*\gamma_d=0$.
By \Cref{lem:odd-sphere-null}, some $D_sa|_F$ is nullhomotopic.
Apply \Cref{top:neighborhood-extension} to this nullhomotopy and
$f=D_sa$ on $Y$.  Its endpoint is the required $h$.  Finally,
$y_\infty\in F$, and every point of $F$ maps to $b$, which
gives the asserted preimage inclusion.
\end{proof}

\begin{remark}[The source of a restricted test]
\label{top:proper-test-interface}
Let $t:Y\to X$ be continuous with $X$ a compact Hausdorff test
space, and write
\[
 f=tqc:S^d\to X,\qquad b=t(y_\infty),
\]
as in \Cref{fig:local-test}. The source over $X\setminus\{b\}$ is precisely
$q^{-1}t^{-1}(X\setminus\{b\})\subset W$, with its given
orientation and $G$-action.  The restriction of $f$ to this
source is proper, since it is the base change of the proper
map $f$ to an open subset of $X$.  Thus its oriented
direct-image pairing is the restriction of the pairing pushed
from $S^d$.  These are the manifold sources used for duality
in the next section; $Y$ enters through ordinary sheaf direct
image.  In \Cref{prop:move-test}, this source lies over the
finite quotient sheets of \Cref{lem:finite-sheets} when $V$ is
chosen there.  The positive degree $s$ may depend on $V$ and
$k$.
\end{remark}
